\documentclass[11pt]{article}
\usepackage[margin=1in]{geometry}
\usepackage{amsmath,amssymb,amsthm,booktabs,longtable,array,calc}
\usepackage[hidelinks]{hyperref}
\usepackage[T1]{fontenc}
\usepackage[utf8]{inputenc}
\newtheorem{theorem}{Theorem}[section]
\newtheorem{proposition}[theorem]{Proposition}
\newtheorem{lemma}[theorem]{Lemma}
\newtheorem{corollary}[theorem]{Corollary}
\theoremstyle{definition}
\newtheorem{definition}[theorem]{Definition}
\theoremstyle{remark}
\newtheorem*{remark}{Remark}
\providecommand{\tightlist}{\setlength{\itemsep}{0pt}\setlength{\parskip}{0pt}}
\begin{document}
\appendix
\section{Exact labelled symmetry maps}\label{appendix-b.-exact-labelled-symmetry-maps}

This appendix gives the literal semantic verification behind Proposition
4.3. Pivot coordinates are treated only in B.6, after the
row-normal-form pivot has been defined in Section 6.

\subsection{\texorpdfstring{B.1. Complement
\(C\)}{B.1. Complement C}}\label{b.1.-complement-c}

Complement fixes every face and crease while replacing every M/V label
by its opposite. Checkerboard orientation and target classes are
unchanged, so every primitive directed precedence reverses. Literal
reversal of the total order reverses every strict inequality and
preserves strict endpoint alternation. Hence \[
T_C(\Pi)=\operatorname{rev}(\Pi)
\] is an exact labelled state bijection.

\subsection{\texorpdfstring{B.2. Row swap
\(R\)}{B.2. Row swap R}}\label{b.2.-row-swap-r}

The face map \(\phi_R\) exchanges \(A_c\) and \(B_c\), maps \(U_i\) to
\(L_i\) and conversely, and fixes \(H_c\). It flips checkerboard
light/dark orientation while preserving the M/V labels. Seam index, and
therefore the \(E/W\) target class, is unchanged. Thus primitive
precedences reverse and \[
T_R(\Pi)=\operatorname{rev}(\phi_R(\Pi)).
\] Incidence, disjointness, target class, and strict-alternation status
are preserved.

\subsection{\texorpdfstring{B.3. Horizontal reflection
\(H\)}{B.3. Horizontal reflection H}}\label{b.3.-horizontal-reflection-h}

The map \(\phi_H\) sends columns by \(c\mapsto n+1-c\), so \[
H_c\mapsto H_{n+1-c},\qquad
U_i\mapsto U_{n-i},\qquad
L_i\mapsto L_{n-i}.
\] Because \(n\) is even, the checkerboard orientation of every face
flips. Horizontal-crease period parity therefore swaps \(h_0\) and
\(h_1\). By contrast, \(n-i\) has the same parity as \(i\), so the \(u\)
and \(l\) period slots and the \(E/W\) target classes are preserved. The
precedence reversal caused by the checkerboard flip gives \[
T_H(\Pi)=\operatorname{rev}(\phi_H(\Pi)).
\] A physical seam \(p\) between columns \(p\) and \(p+1\) maps to the
seam between columns \(n-p\) and \(n-p+1\); therefore \[
p\longmapsto n-p.
\]

\subsection{\texorpdfstring{B.4. The compositions \(CR\) and
\(CH\)}{B.4. The compositions CR and CH}}\label{b.4.-the-compositions-cr-and-ch}

For \(CR\), row swap reverses the directed precedences and complement
reverses them again. The two effects cancel, hence \[
T_{CR}(\Pi)=\phi_R(\Pi).
\] For \(CH\), horizontal reflection and complement likewise contribute
two reversals, so \[
T_{CH}(\Pi)=\phi_H(\Pi).
\] In both cases the same incidence, target-class, and
Butterfly-preservation arguments apply.

\subsection{B.5. Exact two-direction state-set
equivalence}\label{b.5.-exact-two-direction-state-set-equivalence}

Each face transformation is bijective and preserves crease incidence and
disjointness. Corresponding creases have the same folded target class.
Bijective relabelling preserves endpoint positions up to renaming, and
literal order reversal is an order anti-isomorphism; both preserve
whether two disjoint same-target crease intervals strictly alternate.
Thus every semantic clause holds for \(\Pi\) exactly when it holds for
\(T_g(\Pi)\). Since every listed transformation is involutive, \[
\Pi\in F(P,k)\iff T_g(\Pi)\in F(gP,k),
\] and therefore \[
F(gP,k)=T_g(F(P,k)).
\] This is equality of literal labelled state sets, not a symmetry
quotient.

\subsection{B.6. Pivot coordinates after the row-normal-form
definition}\label{b.6.-pivot-coordinates-after-the-row-normal-form-definition}

Let \(R^\sigma(k,p)\) be the row-local normal form of Proposition 6.1,
where \(p\) is the unique physical seam joining the extreme faces of the
row order. Then \[
T_C(R^\sigma(k,p))=R^{-\sigma}(k,p),
\] so complement leaves the numerical pivot unchanged.

Row swap exchanges physical rows without reflecting seam indices. Thus
on an ordered upper/lower pair, \[
(a,b)\longmapsto(b,a)
\] under both \(R\) and \(CR\).

Horizontal reflection sends the physical seam \(p\) to \(n-p\). The row
normal forms satisfy \[
T_H(R^\sigma(k,p))=R^\sigma(k,n-p),
\] \[
T_{CH}(R^\sigma(k,p))=R^{-\sigma}(k,n-p),
\] and therefore \[
(a,b)\longmapsto(n-a,n-b)
\] under \(H\) and \(CH\). These identities transport row-local pivots;
they do not assert that an arbitrary local pivot is globally realizable.

\end{document}
